% GENERATED from paper/draft/02_related_work.md by build.py — do not edit
\section{Background and related work}

\textbf{Delegation and trust in multi-agent systems.} \citet{castelfranchi1998delegation} define delegation as one agent
relying on another for a task, and grade it from weak to strict and from open to closed. Their socio-cognitive account
makes trust a decision to delegate, indexed by trustor, trustee, context, task and goal, and grounded in beliefs about
competence, willingness and opportunity (\citealp{castelfranchi2000trust}). \citet{scerri2002adjustable} cast adjustable
autonomy as a cost-sensitive transfer of control between agent and human, but their model \emph{assumes} a given quality
of agent decisions. \citet{lee2004trust} ask for reliance calibrated to the system's real capability. Work on trust
across tasks shows that trust can be transferred by matching task requirements to agent features
(\citealp{falcone2008generalizing}; \citealp{venanzi2011facing}), that category priors should give way to direct evidence
(\citealp{burnett2010bootstrapping}), and that transfer follows task similarity, from harder tasks to easier ones
(\citealp{soh2020multitask}). \citet{tomasev2026delegation} bring this tradition to LLM agents. They propose authority
graduated by trust, with ``algorithmic circuit breakers'' for revocation and behavioural metrics as trust signals, and
they note that current systems ``lack robust mechanisms for the dynamic assessment of competence''. Their framework is
conceptual, without a measurement procedure or evaluation. \emph{We give an operational instance, tested in a
prospectively specified pilot: competence is estimated per decision class from intervention tests rather than assumed
(Scerri) or read from reputation.}

\textbf{Reasons-responsiveness and meaningful human control.} \citet{fischer1998responsibility} ground responsibility in
a decision mechanism that is moderately reasons-responsive: it recognises reasons and reacts to them across a range of
circumstances. \citet{santonidesio2018mhc} recast this as the \emph{tracking} condition of meaningful human control. A
system must ``demonstrably and verifiably'' respond to the relevant human reasons, and good results can come from
``tracking the wrong property''. \citet{mecacci2020mhc} operationalise tracking for dual-mode vehicles, and
\citet{veluwenkamp2022reasons} argues that the reasons tracked must be normative. CARE-Drive (\citealp{suryana2026caredrive})
is the first behavioural test of tracking for a foundation model. It varies decision facts and prompt-supplied reasons
for a single-step driving model and estimates the sensitivity of its decisions. It has no per-case ground truth for
which fact should decide, no escalation option, and no link to delegation, and it does not separate recognising a reason
from acting on it. \emph{DSA carries the tracking test to multi-step agents and to external evidence whose decisive role
and redundancy are set by construction, relative to one reading of the policy. It is designed to score receptivity and
reactivity separately and to allocate autonomy from the resulting profile; the pilot measures reactivity and reports
receptivity only through the declared-grounds test (§3.4).}

\textbf{Traces, provenance and auditability.} Structured provenance has a long history (\citealp{w3c2013provo}). For LLM
agents, LEDGER builds layered claim-to-evidence graphs over observed sessions and states the problem crisply:
``Visibility, however, is not the same as auditability'' (\citealp{kim2026ledger}). TRACE specifies a typed record that must
exist before any durable state change, and leaves to each consumer the policy that maps a verdict to an action
(\citealp{chang2026trace}). DNative-Twin takes the committed decision as its unit of replay and comparison, and treats
authority as ``granted in advance'' (\citealp{pang2026dnative}). DEMM grades whether a decision's evidence suffices for an
external audit and warns against the ``container fallacy'' (\citealp{solozobov2026demm}). Graph of Trace records the agent's
own declared dependencies as justification edges (\citealp{gao2026graphoftrace}). A recent survey organises the field around
six trust functions: verification, attribution, debugging, safety enforcement, audit and recovery. It lists an
authorization relation as open (\citealp{wang2026survey}). \emph{These systems make decisions inspectable. None tests whether
the recorded grounds are what the decision depends on, and none decides what an agent may do next.}

\textbf{Faithfulness and evidence interventions.} Stated reasons can be unfaithful (\citealp{turpin2023unfaithful};
\citealp{lanham2023measuring}), and interventions are the standard test. \citet{atanasova2023faithfulness} insert inputs and
reconstruct inputs from explanations. \citet{matton2025walk} compare the causal effect of input concepts with how often
explanations cite them. \citet{pawar2026necessary} test necessity and sufficiency of the factors a model cites.
\citet{khanzadeh2026ariadne} edits the agent's own reasoning steps. Closer to agents, AgentAbstain pairs should-act and
should-abstain tasks that differ by one perturbation (\citealp{liu2026agentabstain}). \citet{liao2026provenance} changes the
source authority of evidence in next-action tasks. ``Copies or Sources?'' (\citealp{gao2026copies}) and GraphEcho
(\citealp{wang2026graphecho}) contrast copied with independent sources, and Nautil removes an investigator's cited grounds
against a matched control (\citealp{bi2026nautil}). AUDITA formalises overdetermination with actual causation and shows that
single-cause deletion fails under redundancy (\citealp{du2026audita}). The design itself, an input change with an expected
response, has a longer lineage: metamorphic testing in software engineering (\citealp{segura2016metamorphic}), behavioural
test suites such as CheckList (\citealp{ribeiro2020checklist}), contrast sets (\citealp{gardner2020contrast}) and counterfactually
augmented data (\citealp{kaushik2020cad}), and the diagnosis of models that are right for the wrong reasons through heuristics
(\citealp{mccoy2019right}) or shortcuts (\citealp{geirhos2020shortcut}). \emph{The operators exist. We compose them, adding an
independent source and then removing the original, to separate redundancy from non-use for one decision. We read the
result against the agent-declared grounds and use it prospectively to allocate autonomy, not retrospectively to
attribute blame.}

\textbf{The outcome--process gap.} Several 2026 benchmarks show that outcome checks pass agents whose process is
defective: CXBench in regulated customer-experience workflows (\citealp{yu2026cxbench}), procedure-aware evaluation of τ-bench
(\citealp{cao2026pae}), deterministic evidence-path checks in GroundEval (\citealp{flynt2026groundeval}), and ``lucky passes'' in
software-engineering agents (\citealp{sahoo2026agentlens}). ToolChain-CRC (conformal risk control) proves that calibration on
final answers cannot control unsupported intermediate steps (\citealp{opoku2026toolchain}). \citet{cao2026loopholes} adds a
caution that applies to any benchmark of this kind, ours included: when a natural-language policy is silent, ambiguous
or contradictory, a single gold answer scores policy ambiguity as agent error. \emph{These audits check procedure or
evidence paths. We check whether the action depends on the evidence that should warrant it, which a procedurally
compliant trace can hide, and we show that the same kind of ambiguity can decide which delegation rule appears to work
(§6.1).}

\textbf{Graduated autonomy and certified deployment.} Several systems already widen or narrow agent autonomy from
experience. The Digital Apprentice promotes per-skill tiers when correction rates fall (\citealp{weber2026apprentice}). Hedwig
learns check-in behaviour from a developer's approvals (\citealp{shukla2026hedwig}). Sovereign-OS moves permissions with
audited outcomes (\citealp{yuan2026sovereign}). Minimal Oversight caps authority by a raw success rate, after correcting for
masking (\citealp{azevedo2026minimal}). Trust Before Autonomy prescribes a gated, revocable ladder (\citealp{angelidis2026tba}).
Earned Authority shows that self-produced evidence must not raise a fixed ceiling (\citealp{zhang2026earned}). The Agent
Control Protocol accumulates only penalties from history (\citealp{fernandez2026acp}). Agent Reliability Profiles define
action classes and autonomy tiers for financial services (\citealp{hsu2026arp}). On the statistical side, READY qualifies
confidence-threshold escalation policies on held-out cases, with explicit review cost and an LLM-judged process defect
(\citealp{chatrath2026ready}). \citet{shen2026provable} certifies per model--task pair; Role-Stratified CRC certifies per
stratum (\citealp{rahman2026rolestratified}); Conformal Selective Acting certifies anytime (\citealp{khosravi2026csa}).
\citet{pandey2026certified} gate actions on agreement among sources from separate control domains. \emph{Each delegates
on outcomes, corrections, approvals, confidence, judged traces or source counts. We delegate on interventionally
measured DSA, and we compare against these rules run through one certification engine.}

\textbf{Answer keys as part of the instrument.} Every test of this kind needs an oracle that says what the correct
response is, and the oracle problem is a central difficulty of software testing (\citealp{barr2015oracle}).
\citet{knight1986independence} showed that independently written versions of a program fail on the same inputs far more
often than independence would predict, because their authors find the same cases hard; the lesson carries to two oracles
written from one specification. In machine learning, label errors in test sets change the ranking of models
(\citealp{northcutt2021label}), and disagreement among human annotators is a property of the task rather than noise to be
averaged away (\citealp{plank2022variation}). \citet{jacobs2021measurement} read such problems through measurement theory: a
score is valid only if the construct it operationalises and the procedure that produces its labels are both examined.
Multiverse analysis reports a result under every defensible analytic choice instead of one (\citealp{steegen2016multiverse}),
and formalising law as code shows that statute text leaves choices the formaliser must make explicit
(\citealp{merigoux2021catala}). \emph{We apply these lessons to the normative map of an agent evaluation: we log the readings
of the policy that the oracles encode, report the primary result under every logged combination, and show that the
choice of reading decided which delegation rule appeared to work (§6.1, §8).}


\begin{table*}
\caption{Comparison with the closest work (✓ present; ◐ partial, see the notes; -- absent)}
\small
\begin{tabular}{@{}
  >{\raggedright\arraybackslash}p{(\textwidth - 18\tabcolsep) * \real{0.4646}}
  >{\raggedright\arraybackslash}p{(\textwidth - 18\tabcolsep) * \real{0.0595}}
  >{\raggedright\arraybackslash}p{(\textwidth - 18\tabcolsep) * \real{0.0595}}
  >{\raggedright\arraybackslash}p{(\textwidth - 18\tabcolsep) * \real{0.0595}}
  >{\raggedright\arraybackslash}p{(\textwidth - 18\tabcolsep) * \real{0.0595}}
  >{\raggedright\arraybackslash}p{(\textwidth - 18\tabcolsep) * \real{0.0595}}
  >{\raggedright\arraybackslash}p{(\textwidth - 18\tabcolsep) * \real{0.0595}}
  >{\raggedright\arraybackslash}p{(\textwidth - 18\tabcolsep) * \real{0.0595}}
  >{\raggedright\arraybackslash}p{(\textwidth - 18\tabcolsep) * \real{0.0595}}
  >{\raggedright\arraybackslash}p{(\textwidth - 18\tabcolsep) * \real{0.0595}}@{}}
\toprule\noalign{}
\begin{minipage}[b]{\linewidth}\raggedright
Work
\end{minipage} & \begin{minipage}[b]{\linewidth}\raggedright
(a) decision unit
\end{minipage} & \begin{minipage}[b]{\linewidth}\raggedright
(b) evidence interventions
\end{minipage} & \begin{minipage}[b]{\linewidth}\raggedright
(c) redundancy
\end{minipage} & \begin{minipage}[b]{\linewidth}\raggedright
(d) declared grounds
\end{minipage} & \begin{minipage}[b]{\linewidth}\raggedright
(e) per-class autonomy
\end{minipage} & \begin{minipage}[b]{\linewidth}\raggedright
(f) revocation
\end{minipage} & \begin{minipage}[b]{\linewidth}\raggedright
(g) held-out
\end{minipage} & \begin{minipage}[b]{\linewidth}\raggedright
(h) shift
\end{minipage} & \begin{minipage}[b]{\linewidth}\raggedright
(i) empirical
\end{minipage} \\
\midrule\noalign{}

LEDGER (\citealp{kim2026ledger}) & -- & -- & -- & -- & -- & -- & -- & -- & ◐¹ \\
TRACE (\citealp{chang2026trace}) & ◐² & -- & ◐³ & ✓ & -- & -- & -- & -- & -- \\
DNative-Twin (\citealp{pang2026dnative}) & ✓ & ◐⁴ & -- & -- & -- & -- & -- & -- & -- \\
DEMM (\citealp{solozobov2026demm}) & ✓ & -- & -- & -- & -- & -- & -- & -- & -- \\
Project Ariadne (\citealp{khanzadeh2026ariadne}) & -- & -- & -- & ◐⁵ & -- & -- & -- & -- & ◐⁶ \\
\citet{pawar2026necessary} & ◐⁷ & ◐⁸ & -- & ✓ & -- & -- & -- & -- & ✓ \\
AgentAbstain (\citealp{liu2026agentabstain}) & ◐⁷ & ✓ & -- & -- & -- & -- & -- & -- & ✓ \\
AUDITA (\citealp{du2026audita}) & -- & -- & ✓ & -- & -- & -- & -- & -- & ◐⁹ \\
CARE-Drive (\citealp{suryana2026caredrive}) & ✓ & ◐¹⁰ & -- & -- & -- & -- & -- & -- & ✓ \\
CXBench (\citealp{yu2026cxbench}) & -- & -- & -- & -- & -- & -- & -- & -- & ✓ \\
\citet{pandey2026certified} & ✓ & ◐¹¹ & ✓ & -- & -- & -- & -- & -- & ✓ \\
Intelligent AI Delegation (\citealp{tomasev2026delegation}) & -- & -- & -- & -- & ◐¹² & ◐¹² & -- & -- & -- \\
Trust Before Autonomy (\citealp{angelidis2026tba}) & -- & ◐¹³ & -- & -- & ◐¹² & ◐¹² & -- & -- & -- \\
Digital Apprentice (\citealp{weber2026apprentice}) & -- & -- & -- & -- & ✓ & ✓ & -- & -- & -- \\
Hedwig (\citealp{shukla2026hedwig}) & -- & -- & -- & -- & ✓ & ✓ & -- & -- & ◐¹⁴ \\
Minimal Oversight (\citealp{azevedo2026minimal}) & -- & -- & -- & -- & ✓ & ◐¹⁵ & -- & -- & -- \\
READY (\citealp{chatrath2026ready}) & -- & -- & -- & -- & ◐¹⁶ & ◐¹⁷ & ✓ & -- & ✓ \\
Certified deferral (\citealp{shen2026provable}) & ◐¹⁸ & -- & -- & -- & ◐¹⁹ & -- & ✓ & ◐²⁰ & ◐²¹ \\
Role-Stratified CRC (\citealp{rahman2026rolestratified}) & ◐²² & ◐²³ & -- & -- & ◐²⁴ & -- & ✓ & ✓ & ✓ \\
Conformal Selective Acting (\citealp{khosravi2026csa}) & ◐²⁵ & -- & -- & -- & -- & ◐²⁶ & ◐²⁷ & ✓ & ◐²⁸ \\
\textbf{WorkTrace / DSA (this paper)} & ✓ & ✓ & ✓ & ◐³⁰ & ◐³¹ & ◐²⁹ & ✓ & ✓ & ✓ \\
\bottomrule
\end{tabular}
\vspace{3pt}\par\begingroup\footnotesize Notes to Table 1. Columns: (a) decision unit: a single decision is the unit of analysis; (b) external-evidence
interventions: the agent's external evidence is changed and the agent is re-run; (c) redundancy: redundant or
overdetermined evidence is handled; (d) declared grounds: the agent's own stated grounds are scored; (e) per-class
autonomy: autonomy is allocated per class of decisions or tasks; (f) revocation: autonomy is withdrawn; (g) held-out:
evaluation on held-out cases; (h) shift: evaluation under distribution shift; (i) empirical: empirical study with LLM
agents.\par Footnotes: 1. LEDGER: two qualitative case studies, no metrics. 2. TRACE: one record per claim or commitment. 3. TRACE:
argues that correlated sources add little (effective panel size); no tests. 4. DNative-Twin: replay under declared
changes to recorded state; no LLM agent is re-run. 5. Ariadne: intervenes on the agent's generated reasoning steps, not
on stated grounds about evidence. 6. Ariadne: single-model QA without tools; the report gives two sample sizes (30 and
500) and two headline values for the same rate, so we read it as preliminary. 7. Pawar, AgentAbstain: single decision or
task pair, not multi-step episodes. 8. Pawar: changes input factors of tool-free decisions. 9. AUDITA: live LLM
pipelines, but evidence partly simulated. 10. CARE-Drive: varies facts written into a single prompt. 11. Pandey et al.:
corrupted sources in synthetic payment episodes. 12. Prescribed, not implemented or evaluated. 13. Trust Before
Autonomy: ``unanswerable variants'' correspond to removing the sole decisive source. 14. Hedwig: evaluated on 11
operations with an LLM judge. 15. Minimal Oversight: implicit, through a decaying success-rate estimate; simulation
only. 16. READY: one global threshold per agent and workflow, with per-case accept or escalate. 17. READY:
re-qualification on version changes, no statistical monitoring. 18. Shen: single multiple-choice answers, not agent
decisions. 19. Shen: one certified confidence threshold per model--task pair. 20. Shen: threshold transfer between two
tasks only; the certificate assumes no shift. 21. Shen: eleven small LLMs on multiple-choice QA, without tools. 22.
Role-Stratified CRC: individual argument fields of tool calls. 23. Role-Stratified CRC: prompt injections in tool
outputs, contrasted with clean runs on recorded traces; agents are not re-run. 24. Role-Stratified CRC: thresholds per
argument role, gating field values at run time, not autonomy; recalibration, no revocation. 25. Conformal Selective
Acting: act or abstain on single outputs in one stream. 26. Conformal Selective Acting: all certifications reset at
fixed epoch boundaries; no data-driven revocation. 27. Conformal Selective Acting: risk is controlled online on the same
stream; a held-out split only calibrates the score. 28. Conformal Selective Acting: RLVR-trained LLMs on benchmark
streams, including live online fine-tuning, not agents with tools. 29. This paper: revocation is specified (§4.6) but
not exercised in the pilot. 30. This paper: declared grounds are scored by an LLM judge for factual decisions only
(§6.2); the joint receptivity × reactivity table of §3.5 is not reported. 31. This paper: the pilot exercises granting
only; narrowing and revocation are specified (§4.5, §4.6) but not exercised.\par\endgroup
\end{table*}

